\documentclass[a4paper,11pt]{letter}
\usepackage{siunitx,hyperref,mathpazo,setspace,indentfirst,booktabs}
\usepackage[margin=20mm]{geometry}
\usepackage{newfloat}
\DeclareFloatingEnvironment[fileext=lot]{table}
\hypersetup{colorlinks}
\begin{document}\onehalfspacing
% \pagestyle{headings}
\begin{letter}{}
    \address{IRIS Adlershof\\Humboldt-Universität zu Berlin\\Zum Großen Windkanal 2\\12489 Berlin, Germany}
    
    \opening{Dear Editor,}
    
    I am writing to submit the manuscript entitled `\textit{The Balloon-v1 Model for Metals With Improved Cyclic Response}' for consideration for publication in your journal.
    
    Accurate prediction of cyclic material behaviour is fundamental to various engineering applications.
    Despite decades of research, capturing the full spectrum of cyclic phenomena, including cyclic hardening/softening, ratcheting, and stress relaxation, within a single, versatile constitutive framework remains a significant challenge in computational mechanics.
    
    The manuscript attempts to address this challenge by proposing a novel constitutive model that builds upon the subloading surface theory.
    The key innovation lies in a new isotropic hardening formulation that leverages loading reversal information embedded in the normal yield ratio offered by the subloading surface theory.
    This approach provides an alternative (simpler and better) approach to capture complex cyclic responses compared to existing approaches.
    
    The proposed Balloon-v1 model offers:
    \begin{enumerate}
        \item a unified framework for diverse cyclic phenomena without requiring separate formulations for different loading scenarios,
        \item a physically motivated approach based on revised subloading surface theory, and
        \item validation against experimental data demonstrating its versatility and accuracy.
    \end{enumerate}
    
    This work aligns well with the scope of \texttt{Acta Mechanica} in computational mechanics and constitutive modelling.
    The combination of theoretical development, numerical implementation, and experimental validation makes it suitable for your readership of researchers and practitioners in plasticity and constitutive modelling for metallic materials.
    
    The manuscript is a continuation of the previous work \href{https://link.springer.com/article/10.1007/s00707-025-04339-0}{10.1007/s00707-025-04339-0}.
    Please feel free to appoint reviewers as there are no conflicts of interest.
    Reviewers who reviewed the previous manuscript are preferable since they are familiar with the background of this work.    
    Thank you for considering our manuscript.
    
    \signature{Dr. Theodore Chang}
    
    \closing{Warm regards,}
\end{letter}
\clearpage
\end{document}
